While the Reeb graph construction is the original conceptual foundation for the Mapper 
algorithm, much of the Mapper construction's theoretical and geometric intuition 
can be derived from \emph{Leray’s Nerve Theorem}. In particular, Mapper can be viewed 
as an attempt to reconstruct, from sampled data, a good cover whose nerve 
approximates the underlying manifold $X$. Each choice of lens function, cover, 
and clustering method defines how this discrete cover is built. And since 
the cover is assumed to be \emph{good} the choices of hyperparameters induce priors on the 
connectivity structure for the underlying manifold mapper approximates.

The lens function $f : D \rightarrow B$ determines how the dataset is projected 
into a base space $B$, effectively specifying the perspective through which 
the topology of $X$ is observed. Depending on the choice of $f$
the same dataset $D$ can produce Mapper complexes highlighting different features 
of $X$. Intuitively, $f$ defines which directions in data space reveal relevant 
geometric or functional variation. A well-chosen filter preserves separations 
inherent to $X$, while a poorly chosen one may cause topologically distinct 
regions to overlap in $B$, introducing artificial simplicies in the resulting Mapper 
complex.

$$\text{[include dataset example]}$$

The covering scheme $\mathcal{C}$ on the image $f[D]$ controls both 
the resolution of the topological reconstruction. Generally speaking, a coarse cover 
with few overlapping intervals yields a low-resolution summary that can smooth 
out smaller-scale structures, while a fine cover with high overlap can expose more 
localized features but increases sensitivity to the sampling noise. 
Through the preimages $f^{-1}(U_i)$, the cover induces a family of overlapping 
subsets of $D$, forming what we call the \emph{cluster cover} of $X$. G-Mapper 
demonstrates this by adaptively refining the cover in certain regions assuming 
noise is normally distributed.

$$\text{[include example]}$$

The clustering method serves as a discrete approximation to 
path-connected components within each preimage $f^{-1}(U_i)$. This choice 
directly influences how Mapper interprets local connectivity: 
density-based algorithms such as DBSCAN, HDBSCAN, or single-linkage clustering 
tend to approximate topological connectedness by grouping dense or chain-connected 
regions, while partitional methods such as $K$-Means impose a fixed number of 
clusters per region. The latter approach introduces strong geometric assumptions 
(e.g., roughly convex clusters), which may guarantee that each cluster is 
contractible (helpful for forming a \emph{good cover}) but potentially deviate from 
the manifold’s true topology. Density-based methods, while more faithful to the 
notion of connected components from the Reeb construction, can still fail to yield a good cover due to 
the chaining effect introducing non-contractible intersections between clusters,
or suffer from increased sensitivity to noise and scale. 

$$\text{[include example]}$$

\subsection{Parameter Selection}
The triple $(f, \mathcal{C}, \text{cluster})$ defines a discrete 
approximation to the topology of $X$ from a sampling $D$. If the cluster cover forms a \emph{good} cover, then 
the Mapper complex is homotopy equivalent to $X$. In practice, this equivalence 
is approximate: the extent to which Mapper recovers the true topology depends 
on how the lens function partitions the data, how the cover refines the image, 
and how the clustering algorithm captures local connectivity. Consequently, 
the Mapper construction can be viewed as exploring a family of topological 
summaries parameterized by the hyperparameter space $\mathcal{H}$, where each 
choice of parameters provides a different resolution and perspective on the 
shape of the underlying space $X$. Typically this does not yield a filtration...
